\section*{Aufgabenstellung}
Stellen Sie sicher, dass Sie alle Punkte in der Aufgabenstellung berücksichtigt haben (Upload Dateiname, Prüfsummenvergleich (md5) etc.
Welche Sleuthkit und xmount Version kommt bei Ihnen zum Einsatz?

\section*{Lösung}

\subsection*{MD5 Checksumme}

Berechnete Checksumme der Images mithilfe von md5sum:
\begin{description}
	\item[./images/klausur/20260328\_apl.E01] 064832e41e7a5c6de39cf78673cd7083
	\item[./images/klausur/20260328\_apl.E02] a84c87d1e70707e15ee55693aacf189c
	\item[./images/klausur/20260328\_apl.E03] d8ebfe3194db454f7802f5138d7eb396
	\item[./images/klausur/20260328\_apl.E04] db2a7cde242f3a53c3cfa9a0b2bc8de9
	\item[./images/klausur/20260328\_apl.E05] d31f769e0d8daebf7b55e4d083d4dab6
	\item[./images/klausur/20260328\_apl.E06] 6ef5a2879cc4e08e0393bd9152f91378
	\item[./images/klausur/20260328\_apl.E07] 8d3cd4639c9b8107fda96f2eec5b1c8f
	\item[./images/klausur/20260328\_server.E01] bf18a6bdc9e56b1b30334800b254af15
\end{description}

\subsection*{Sleuthkit Version}
\begin{lstlisting}[language=bash]
root@Alpha-17:/home/luis/Programming/GitHub/University/DigitalForensic#   mmls -V
The Sleuth Kit ver 4.12.1
\end{lstlisting}

\subsection*{xmount Version}
\begin{lstlisting}[language=bash]
root@Alpha-17:/home/luis/Programming/GitHub/University/DigitalForensic# xmount --version
xmount v0.7.6 Copyright (c) 2008-2018 by Gillen Daniel <gillen.dan@pinguin.lu>

  compile timestamp: Apr  1 2024 08:38:03
  gcc version: 13.2.0
  loaded input libraries:
    - libxmount_input_aff.so supporting "aff"
    - libxmount_input_ewf.so supporting "ewf"
    - libxmount_input_aewf.so supporting "aewf"
    - libxmount_input_raw.so supporting "raw", "dd"
    - libxmount_input_aaff.so supporting "aaff"
  loaded morphing libraries:
    - libxmount_morphing_combine.so supporting "combine"
    - libxmount_morphing_raid.so supporting "raid0"
    - libxmount_morphing_unallocated.so supporting "unallocated"
\end{lstlisting}

\subsection*{binwalk Version}
Zusätzlich habe ich die binwalk Version überprüft, da diese für die Analyse der Images relevant sein könnte:
\begin{lstlisting}[language=bash]
luis@Alpha-17:~/Programming/GitHub/University/DigitalForensic$ pip show binwalk 2>/dev/null | grep -i version
Version: 2.3.3
\end{lstlisting}
